\documentclass{amsart}
\usepackage{todonotes}
\usepackage{xypic}
\usepackage{tikz-cd}
\usepackage{graphicx, amsmath, amssymb, amsthm}
%\usepackage{showkeys}

\usepackage{mathrsfs}
\usepackage{cite}
\usepackage{comment}
\usepackage{tikz,pgf}
\usepackage{arcs}
\usepackage{caption}
\usetikzlibrary{positioning}
\usetikzlibrary{arrows}
\usepackage{float}
\usepackage{hyperref} 

\usetikzlibrary{arrows}
\usetikzlibrary{decorations.pathreplacing}
\usetikzlibrary{calc}
\usepackage{graphicx} 

%Basic Notation

\newcommand{\mycases}[1]{\left\{\begin{array}{ll}#1\end{array}\right.}
\newcommand{\Z}{{\mathbb  Z}}
\newcommand{\N}{{\mathbb N}}
\newcommand{\R}{{\mathbb R}}
\newcommand{\F}{{\mathbb F}}
\newcommand{\MUR}{MU_{\R}}
\newcommand{\MSpin}{MSpin}
\newcommand{\Spin}{Spin}
\newcommand{\XiO}{\Xi O}
\newcommand{\smashover}[1]{\underset{#1}{\wedge}}
\newcommand{\timesover}[1]{\underset{#1}{\times}}
\newcommand{\Boxover}[1]{\underset{#1}{\Box}}
\newcommand{\otimesover}[1]{\underset{#1}{\otimes}}

%Functors
%\newcommand{\Ind}{\big\uparrow} %Matching Mackey notation
%\newcommand{\Res}{\big\downarrow} %Matching Mackey notation
%\newcommand{\ind}{\uparrow}
%\newcommand{\res}{\downarrow}
\DeclareMathOperator{\Hom}{Hom}
\DeclareMathOperator{\Der}{Der}
\DeclareMathOperator{\Ext}{Ext}
\DeclareMathOperator{\Tor}{Tor}
\DeclareMathOperator{\coker}{coker}
\DeclareMathOperator{\Map}{Map}
\DeclareMathOperator{\Sym}{Sym}
\DeclareMathOperator{\holim}{holim}
\DeclareMathOperator{\colim}{colim}
\DeclareMathOperator{\Stab}{Stab}
\DeclareMathOperator{\Spec}{Spec}
\DeclareMathOperator{\Tel}{Tel}
\DeclareMathOperator{\Sub}{Sub}

\newcommand*{\RO}{RO}

\newcommand{\norm}{N}
\newcommand{\normR}{{}_R\norm}

%Mackey Functor Structure
\newcommand{\res}{res}
\newcommand{\tr}{tr}
\newcommand{\Tr}{tr}

%Groups
\newcommand{\Cp}[1]{C_{p^{#1}}}
%\newcommand{\Cp}[1]{G_{#1}}
\newcommand{\Cpn}{\Cp{n}}
\newcommand{\cp}[1]{p^{#1}}

%Mackey Functor Names
\newcommand{\m}[1]{{\protect\underline{#1}}}
\newcommand{\mZ}{\m{\Z}}
\newcommand{\mM}{\m{M}}
\newcommand{\mG}{\m{G}}
\newcommand{\mR}{\m{R}}
\newcommand{\mB}{\m{B}}
\newcommand{\mJ}{\m{J}}
\newcommand{\mH}{\m{H}}
\newcommand{\mN}{\m{N}}
\newcommand{\mpi}{\m{\pi}}
\newcommand{\mn}{\m{n}}
\newcommand{\mA}{\m{A}}
\newcommand{\mX}{\m{X}}
\newcommand{\mSpec}{\m{\Spec}}
\newcommand{\mY}{\m{Y}}
\newcommand{\mGamma}{\m{\Gamma}}
\newcommand{\mSet}{\m{\Set}}
\newcommand{\mcC}{\m{\cC}}
\newcommand{\mTr}{\m{\Tr}}
\DeclareMathOperator{\hull}{hull}

%Mathcal Commands
\newcommand{\cc}[1]{\mathcal #1}
\newcommand{\cC}{\cc{C}}
\newcommand{\ccD}{\cc{D}}
\newcommand{\cF}{\cc{F}}
\newcommand{\cO}{\cc{O}}
\newcommand{\ccZ}{\cc{Z}}
\newcommand{\cP}{\cc{P}}
\newcommand{\cA}{\cc{A}}
%\newcommand{\cL}{\cc{L}}

%AJB kludges
\newcommand{\aC}{\cC}
\newcommand{\aD}{\ccD}
\newcommand{\ob}{\textrm{ob}}
\newcommand{\id}{\textrm{id}}

%Category Names
\newcommand{\TC}{\textnormal{TC}}
\newcommand{\THH}{\textnormal{THH}}
\newcommand{\AQ}{\textnormal{AQ}}
\newcommand{\TAQ}{\textnormal{TAQ}}
\newcommand{\CoInd}{\textnormal{CoInd}}
\newcommand{\Ind}{\textnormal{Ind}}
\newcommand{\HH}{\textnormal{HH}}
\newcommand{\GL}{\textnormal{GL}}
\newcommand{\Sp}{\mathcal Sp}
\newcommand{\Emb}{\textnormal{Emb}}
\newcommand{\Set}{\mathcal Set}
\newcommand{\Com}{\mathcal Comm}
\newcommand{\Comm}{\mathcal Comm}
\newcommand{\Nuca}{\mathcal NUCA}
\newcommand{\Ninfty}{N_\infty}
\newcommand{\Coef}{\mathcal Coef}
\newcommand{\Ens}{\mathcal Ens}
\newcommand{\Top}{\mathcal Top}
\newcommand{\Ab}{\mathcal Ab}
\newcommand{\Alg}{\mathcal Alg}
\newcommand{\cOrb}{\mathcal Orb}
\newcommand{\Scheme}{\mathcal Scheme}
\newcommand{\Tamb}{\mathcal Tamb}
\newcommand{\OTamb}{\cO\mhyphen\Tamb}
\newcommand{\Mackey}{\mathcal Mackey}
\newcommand{\Cat}{\mathcal Cat}
\newcommand{\Symm}{\mathcal Sym}
\newcommand{\mRmod}{\m{R}\mhyphen\mathcal Mod}
\newcommand{\GComm}{G\mhyphen\Comm}
\newcommand{\Green}{\mathcal Green}
\newcommand{\mScheme}{\m{\Scheme}}
\newcommand{\Burn}{\mathcal Burn}
\newcommand{\Coeff}{\mathcal Coeff}
\newcommand{\mCoeff}{\m{\Coeff}}
\newcommand{\Mod}{\mathcal Mod}
\newcommand{\Mods}[1]{#1\mhyphen\Mod}
\newcommand{\RMod}{\Mods{R}}
\newcommand{\mRMod}{R\mhyphen\m{\Mod}}
\newcommand{\STamb}{\m{S}\mhyphen\Tamb}
\newcommand{\SAb}{\m{S}\mhyphen\Ab}
\newcommand{\SMackey}{\m{S}\mhyphen\Mackey}
\newcommand{\Aug}[1]{#1\mhyphen\mathcal Aug}
\newcommand{\OmegaOneO}{\Omega^{1,\mathcal O}_{\m{R}/\m{S}}}
\newcommand{\OmegaOneG}{\Omega^{1,G}_{\m{R}/\m{S}}}
\newcommand{\semidirect}{\ltimes}
\DeclareMathOperator{\conn}{conn}
\DeclareMathOperator{\gconn}{gconn}
\newcommand{\Fin}{\mathcal Fin}
\newcommand*{\IntOp}{\mathcal IntOp}
\newcommand*{\Tran}{\mathcal Tran}
\newcommand*{\SubPo}{\mathcal SubPo}
\DeclareMathOperator{\claw}{claw}
\DeclareMathOperator{\Aut}{Aut}
\DeclareMathOperator{\Cata}{Cat}
\newcommand{\MRC}{MRC}
\newcommand{\MRCcOMackey}{\MRC_{\cO}\mhyphen\Mackey}
\DeclareMathOperator{\End}{End}

\DeclareMathOperator{\mgconn}{g\m{\conn}}

\newcommand{\leftadjoint}{\dashv}
\newcommand{\To}{\Rightarrow}

\newcommand{\AN}[1]{{}_{#1}N}

\newcommand{\defn}[1]{{\underline{#1}}}

\newcommand{\mC}{\m{\cC}}

%Hyphens
\mathchardef\mhyphen=45

 %Random Macros
\newcommand{\EM}{Eilenberg-Mac~Lane}

%Theorem Styles
\numberwithin{equation}{section}

\newtheorem{theorem}{Theorem}[section]
\newtheorem{lemma}[theorem]{Lemma}
\newtheorem{corollary}[theorem]{Corollary}
\newtheorem{cor}[theorem]{Corollary}
\newtheorem{proposition}[theorem]{Proposition}
\newtheorem{conjecture}[theorem]{Conjecture}
\newtheorem{question}[theorem]{Question}

\newtheorem*{theorem*}{Theorem}

\theoremstyle{remark}
\newtheorem{remark}[theorem]{Remark}
\newtheorem{example}[theorem]{Example}
\newtheorem{notation}[theorem]{Notation}
\newtheorem{observation}[theorem]{Observation}
\newtheorem{rant}[theorem]{Rant}
\newtheorem{warning}[theorem]{Warning}


\theoremstyle{definition}
\newtheorem{assumption}[theorem]{Assumption}
\newtheorem{axiom}[theorem]{Axiom}
\newtheorem{definition}[theorem]{Definition}


%Miscellaneous macros
\newcommand{\SESM}{Eilenberg--Mac~Lane}
\newcommand{\eqvr}{equivariant}
\newcommand{\defemph}[1]{\textbf{#1}} 
\newcommand{\Ho}{\textrm{Ho}}

\newcommand{\sma}{\otimes}
\newcommand{\sP}{\mathcal{P}}
\newcommand{\htp}{\simeq}
\newcommand{\aO}{\mathcal{O}}
\newcommand{\aU}{\mathcal{U}}
\newcommand{\bE}{\mathbb{E}}
\newcommand{\bO}{\mathbb{O}}
\newcommand{\bU}{\mathbb{U}}
\newcommand{\GTop}{\textrm{GTop}}
\newcommand{\bN}{\mathbb{N}}
\newcommand{\Ev}{\textrm{Ev}}

\title{Generalized equivariant slice categories}

\author{Andrew J.~Blumberg}
\author{Michael A.~Hill}
\author{Tyler Lawson}

\begin{document}

\section{Indexed slice categories}\label{sec: Indexed Slice
  Categories}

(Excerpt from ``Generalized equivariant slice categories'', with Mike
Hill and Tyler Lawson.)

\subsection{Transfer and indexing systems}

We begin with an ahistorical but geodesic summary of transfer systems
and indexing systems.

\begin{definition}[{\cite{BBR}, \cite{Rubin21}}]
    A {\emph{transfer system}} on \(G\) is a partial order we will denote by \(\to\) on \(\Sub(G)\) satisfying three properties:
    \begin{enumerate}
        \item it refines subgroup inclusion: if \(H\to K\), then \(H\subseteq K\),
        \item it is conjugation invariant: if \(H\to K\) and \(g\in G\), then \(gHg^{-1}\to gKg^{-1}\), and
        \item it is closed under restriction: if \(H\to K\) and \(J\subseteq K\), then \(H\cap J\to J\).
    \end{enumerate}
\end{definition}

The collection of all transfer systems on \(G\) forms a poset under refinement, and we will use \(\leq\) for the partial order here.

\begin{definition}
    Let \(\cO\) be a transfer system on \(G\). A finite \(H\)-set 
    \[
        T=\coprod_{i} H/K_i
    \]
    is {\emph{admissible}} for \(\cO\) if for all \(i\), \(K_i\to H\). The collection of admissible \(H\)-sets for \(\cO\) will be denoted \(\cO(H)\). The collection of all \(\cO(H)\) as \(H\) varies gives an {\emph{indexing system}}.
\end{definition}

The admissible sets of $\cO$ are closely connected to the norms
structured by an \(\Ninfty\) operad; we will usually also abusively
denote the operad by \(\cO\). Here $i^H_* \colon \Sp^G \to \Sp^H$
denotes the pullback functor along the inclusion $H \to G$ and $N_H^G
\colon \Sp^H \to \Sp^G$ denotes the Hill-Hopkins-Ravenel
norm~\cite{HHR}.

\begin{definition}
    For a finite \(G\)-set \(T\), we define the \(T\)-norm
    \[
        N^T\colon\Sp^G\to\Sp^G
    \]
    inductively by the formulas
    \begin{enumerate}
        \item \(N^{G/H}(E)=N_H^G i_H^\ast(E)\), and
        \item \(N^{T_0\amalg T_1}(E)=N^{T_0}(E)\otimes N^{T_1}(E)\).
    \end{enumerate}
\end{definition}

\subsection{\texorpdfstring{\(\cO\)}{O}-slice filtration}

We now define the slice filtration relative to an indexing system
$\cO$.  We are going to use equivariant localization (more
specifically, nullification) to construct the relative slice towers.
Recall that in the equivariant context, we define local and acylic
objects in terms of conditions on the $G$-space of maps rather than
the non-equivariant space of maps.  The acylic objects form an
equivariant localizing subcategory.  Recall that given a set of
objects in $\Sp^G$, we define the equivariant localizing subcategory
generated by these objects to be the full subcategory of $\Sp^G$
constructed as the closure under homotopy colimits, retracts, and
tensors with orbit spectra.

\begin{definition}
    If \(\cO\) is an indexing system, then let \(\tau_{\geq n}^{\cO}\) be the equivariant localizing subcategory of $\Sp^G$ generated by 
    \[
    \Big\{ G_+\otimesover{H} N^T S^1\mid T\in \cO(H), |T|\geq n\Big\}.
    \]
    This is the category of \emph{\(\cO\)-slice \(n\)-connective spectra}.
\end{definition}

\begin{remark}
    Given a finite \(G\)-set \(T\), we have an equivariant homeomorphism
    \[
        N^T S^1\cong S^{\mathbb R\cdot T},
    \]
    the representation sphere associated to the permutation representation of \(T\). This means that the \(\cO\)-slice \(n\)-connective spectra can be equivalently viewed as being generated by the representation spheres associated to the permutation representations for admissible sets of cardinality at least \(n\). 
\end{remark}

Viewing this instead as a diagram of localizing subcategories (i.e., as a categorical Mackey functor), we are forming the equivariant localizing subcategory generated at \(G/H\) by \(N^T S^1\) for all admissible \(H\)-sets \(T\) of cardinality at least \(n\).

For the next definition, recall that the nullification at a set of objects $\{S_i\}$ in $\Sp^G$ is the left Bousfield localization at the set of terminal maps $\{S_i \to \ast\}$.

\begin{definition}
    If \(\cO\) is an indexing system, then:
    \begin{itemize}

    \item The \(n\)th \(\cO\)-slice truncation is the functor
    \[
        P^n_{\cO}\colon \Sp^G_{\geq 0}\to\Sp^G_{\geq 0}
    \]
    that is the nullification killing \(\tau_{\geq (n+1)}^{\cO}\). 

\item The \(n\)th \(\cO\)-slice cover is the functor
\[
P_n^{\cO}\colon \Sp^G_{\geq 0}\to\Sp^G_{\geq 0}
\]
defined to be the (homotopy) fiber of the natural map \(Id\Rightarrow P^{n-1}_{\cO}\). 
\end{itemize}
\end{definition}

The truncation functors are related in the evident fashion as $n$ varies.

\begin{proposition}
    For each \(n\geq 0\), we have a natural transformation 
    \[
    P^n_{\cO}(\mhyphen)\Rightarrow P^{n-1}_{\cO}(\mhyphen).
    \]
    These are compatible with the natural nullification functors 
    \[
    Id\Rightarrow P^n_{\cO}(\mhyphen).
    \]
    
    For a connective $G$-spectrum $E$, the natural map 
    \[
        E\to \lim_{\longleftarrow} P^n_{\cO}(E)
    \]
    is always a weak equivalence.
\end{proposition}
\begin{proof}
    The inclusion of categories \(\tau_{n+2}^{\cO}\subset\tau_{n+1}^{\cO}\) induces a natural transformation the other way of nullification functors.  Since we can factor the nullification functor \(P^n_{\cO}\) via this inclusion, the first two statements follow.
    
    For the second, we note that the Postnikov connectivity of \(G_+\otimesover{H} N^T S^1\) for a finite \(H\)-set $T$ is \(|T/H|\). As \(n\) goes to infinity, this also does (at worst as \(|T|/|H|)\). In particular, the map
    \[
        E\to P^n_\cO(E)
    \]
    has coconnectivity going to infinity.
\end{proof}

For any bounded below spectrum \(K\), the same argument shows that the natural map
\[
    K\sma E\to \lim_{\longleftarrow}\big(K\sma P^{n}E\big)
\]
is an equivalence. 

\begin{definition}
    A \(G\)-spectrum \(E\) is an \(\cO\)-\(n\)-slice if 
    \begin{enumerate}
        \item it is in \(\tau_{\geq n}^{\cO}\), and 
        \item the natural map
    \[
    E\to P^n_{\cO}E
    \]
    is an equivalence.
    \end{enumerate}
\end{definition}

\begin{proposition}
For any indexing system \(\cO\), the ordinary suspension yields maps
    \[
        \Sigma\colon\tau_{\geq k}^{\cO}\to\tau_{\geq (k+1)}^{\cO}.
    \]
\end{proposition}
\begin{proof}
    Since suspension commutes with homotopy colimits and induction, it suffices to show this on the generators \(N^T S^1\) as $T$ varies over the admissible sets of $\cO$.  Since $\Sigma N^{T} S^1 \htp N^{T\amalg \ast}S^1$, the result follows: if $T$ is admissible and of cardinality at least \(k\) then $T \amalg \ast$ is admissible and has cardinality at least $k+1$.
\end{proof}

\begin{corollary}
    For any \(k\geq 0\), the \(\infty\)-category of \(\cO\)-\(k\)-slices is discrete.
\end{corollary}
\begin{proof}
    If \(E,E'\) are \(\cO\)-\(k\)-slices, then they are both in \(\tau_{\geq k}^{\cO}\). By the usual adjunctions, for all \(n\geq 1\), the higher homotopy group \(\pi_n\) of the mapping space are given by
    \[
        \pi_n\Map(E,E')=[\Sigma^n E,E']^G=0,
    \]
    since the preceding proposition implies that \(\Sigma^n E\in\tau_{\geq (k+n)}^{\cO}\).
\end{proof}

\begin{definition}
We define {\emph{\(n\)\textsuperscript{th} \(\cO\)-slice}} of a connective \(G\)-spectrum \(E\), denoted \(P^n_{n,\cO}(E)\), to be the homotopy fiber of the natural map
\[
P^n_{\cO}(E)\to P^{n-1}_{\cO}(E).
\]
\end{definition}

\section{Characterizing slice towers via connectivity}\label{sec: connectivity}

\subsection{Geometric fixed points and slice connectivity}

We can detect slice connectivity in terms of the connectivity of the geometric fixed points~\cite{HillYarnall, WilsonSlice}.  To express this, it is convenient to define the following function capturing the structure of the indexing system.

\begin{definition}
    For any transfer system \(\cO\), we define the {\em characteristic function} of $\cO$ 
    \[
        \chi^{\cO}\colon\Sub(G)\to\Sub(G)
    \]
    by the formula
    \[
        \chi^{\cO}(H)=\min\{K\mid K\to H\}=\bigcap_{K\to H} K.
    \]
\end{definition}

\subsubsection{The geometric fixed points of \(\tau_{\geq n}^{\cO}\)}

Stable equivalences in $\Sp^G$ can be detected as maps that induce non-equivariant stable equivalences on passage to geometric fixed points for all (closed) subgroups of $G$.  It should thus be very plausible that the connectivity of geometric fixed points is a central notion.

\begin{definition}
    For a \(G\)-spectrum \(E\), let the {\em geometric connectivity}, denoted \(\mgconn(E)\), be the function from subgroups of \(G\) to \(\Z\cup\{\pm\infty\}\) defined by
    \[
        \mgconn(E)(H):= \conn\big(\phi^H(E)\big).
    \]
\end{definition}

\begin{lemma}\label{lem: Slice Connectivity implies Geometric connectivity}
    Let \(\cO\) be a transfer system. If \(E\in\tau_{\geq n}^{\cO}\), then for all \(H\subset G\), 
    \[
        [H:\chi^{\cO}(H)]\cdot \mgconn(E)(H)\geq n.
    \]
\end{lemma}
\begin{proof}
    By restriction, it suffices to show this for \(H=G\). Since the geometric fixed points preserve homotopy colimits and extensions, it suffices to show this for generators.  Next, since geometric fixed points applied to an induced $G$-spectrum vanish, we are reduced to considering the case of \(N^T S^1\) for \(T\) an admissible \(G\)-set of cardinality at least \(n\). Decompose $T$ as
    \[
    T=\sum_{H} n_H G/H.
    \]
    The geometric fixed points of \(N^T S^1\) are \(S^{|T/G|}\), and in this case, we have
    \[
    |T/G|=\sum_H n_H.
    \]
    We have by assumption
    \[
    |T|=\sum_{H} n_H[G:H]\geq n,
    \]
    and by definition, \([G:\chi^{\cO}(H)]\) is the maximal element in 
    \[
    \big\{[G:H]\mid G/H\in\cO(\ast)\big\}
    \]
    (and in fact, all others divide it). This gives inequalities
    \[
    [G:\chi^{\cO}(G)]\cdot \sum_{H} n_H\geq \sum_{H} n_H [G:H]\geq n,
    \]
    as desired.
\end{proof}

\begin{remark}
    If \(\chi^\cO(G)=\{e\}\), then we recover \cite[Theorem 2.5]{HillYarnall}.
\end{remark}

For the converse, we can again use isotropy separation, studying the cofiber sequence
\[
E\mathcal F_+\sma E\to E\to \tilde{E}\mathcal F\sma E.
\]  
The spectrum \(E\mathcal F_+\sma E\) is built out of pieces of the form \(G/H_+\sma E\), so this is in a localizing subcategory if and only if the restrictions are.

\begin{lemma}\label{lem: Localizing Cats and Families}
    Let \(\mathcal F\) be a family, and let \(\tau\) be an equivariant localizing subcategory. If \(E\) is any \(G\)-spectrum such that for all \(H\in\mathcal F\), \(i_H^\ast E\in i_H^{\ast}\tau\), then
    \[
    (E\mathcal F_+\sma E)\in\tau
    \]
\end{lemma}
\begin{proof}
    This follows by the same proof as \cite[Lemma 2.4]{HillYarnall}: the spectrum \(E\mathcal F_+\otimes E\) is in the localizing category generated by \(G/H_+\otimes E\) for \(H\in\mathcal F\). By assumption, we have an inclusion
    \[
        G/H_+\otimes E\cong G_+\otimesover{H} i_H^\ast E\in\tau.
    \]
\end{proof}

\subsubsection*{The \texorpdfstring{\(\cO\)}{O}-slices of geometric spectra}
Our argument will use downward induction on the subgroup lattice, so we will need to understand the \(\cO\)-slice connectivity of \(\tilde{E}\mathcal P\otimes E\), where \(\mathcal P\) is the family of proper subgroups of \(G\).  Recall that a \(G\)-spectrum \(E\) is called ``geometric'' if the natural map
\[
	E\to \tilde{E}\mathcal P\otimes E
\]
is an equivalence \cite[Definition 6.10]{HillSlicePrimer}, and a Mackey functor \(\mM\) is geometric if \(H\mM\) is. The proof of \cite[Theorem 6.7]{HillSlicePrimer} goes through essentially without change to show the following.

\begin{lemma}\label{lem: Slice Connectivity of Geometric Things}%[{\cite[Theorem 6.7]{HillSlicePrimer}}]
    Let \(\mM\) be a geometric Mackey functor. For any \(\cO\), 
    \[
        \Sigma^k H\mM
    \]
    is a \(k\cdot [G:\chi^{\cO}G]\)-\(\cO\)-slice.
\end{lemma}
\begin{proof}
    Since \(\mM\) is geometric, we have that for any finite \(G\)-set \(T\), the natural map
   \[
    S^{|T/G|}\hookrightarrow N^T S^1
   \]
   given by the inclusion of fixed points induces an equivalence
   \[
   S^{|T/G|}\otimes H\mM\to N^TS^1\otimes H\mM.
   \]
   We can bound the \(\cO\)-slice connectivity from below by choosing an \(\cO\)-admissible \(T\) with \(|T|\) as large as possible so that \(|T/G|=k\) is fixed. This is again achieved by taking
   \[
   T=k G/\chi^{\cO}(G),
   \]
   since \(\chi^{\cO}(G)\) is the minimal subgroup \(H\) such that \(H\to G\). This shows us that
   \[
   \Sigma^k H\mM\in\tau_{\geq k[G: \chi^{\cO}(G)]}^{\cO}.
   \]
   
   For the upper bound, consider an admissible \(G\)-set \(T\) such that 
   \[
   |T|>k[G:\chi^{\cO}(G)].
   \]
   Since \(k[G:\chi^{\cO}(G)]\) is the largest cardinality of an admissible \(G\)-set with \(k\)-orbits, we deduce that \(|T/G|>k\). Since \(\mM\) is geometric, we therefore deduce
   \[
   [N^TS^1,\Sigma^kH\mM]^G\cong [\Phi^G N^TS^1,\Sigma^k H\mM(G/G)]
   \cong [S^{|T/G|},\Sigma^{k}H\mM(G/G)]=0.
   \]
   This shows that \(H\mM\) is a \(k[G: \chi^{\cO}(G)]\)-slice.
\end{proof}

\subsection{Rewriting \texorpdfstring{\(\cO\)}{O}-slice connectivity}
Putting these together, we get the full \(\cO\)-slice version of \cite[Theorem 2.5]{HillYarnall}. 

\begin{theorem}\label{thm: GFP Reformulation}
    A \(G\)-spectrum \(E\) is in \(\tau_{\geq n}^{\cO}\) if and only if for all \(H\subset G\),
    \[
    [H:\chi^{\cO}(H)]\cdot\mgconn(E)(H)\geq n.
    \]
\end{theorem}

\begin{proof}
    The proof is essentially that of \cite[Theorem 2.5]{HillYarnall}. The forward direction is  Lemma~\ref{lem: Slice Connectivity implies Geometric connectivity}. 
    
    For the other direction, let \(E\) be a spectrum with the prescribed geometric connectivities. Consider the isotropy separation sequence
    \[
        EP_+\otimes E\to E\to \tilde{E}\mathcal P\otimes E.
    \]
    By Lemma~\ref{lem: Slice Connectivity of Geometric Things}, the \(\cO\)-slice connectivity of \(\tilde{E}\mathcal P\otimes E\) is at least \(n\). By induction on the subgroup lattice, Lemma~\ref{lem: Localizing Cats and Families} shows that \(EP_+\otimes E\) also has \(\cO\)-slice connectivity \(n\). Since localizing categories are closed under extensions, this implies that \(E\) has \(\cO\)-slice connectivity \(n\).   
\end{proof}

Rewriting this slightly, we have a way to describe the slice connectivity of an arbitrary \(0\)-connective spectrum.

\begin{corollary}\label{cor: Determining Slice Connectivity}
    If \(E\in\Sp^{G}_{\geq 0}\), then let 
    \[
        n=\min_{H\subseteq G}\big\{[H:\chi^{\cO}(H)]\cdot\mgconn(E)(H)\big\}.
    \]
    Then \(E\in\tau_{\geq n}^{\cO}\).
\end{corollary}

\begin{thebibliography}{1}

\bibitem{BBR}
S.~Balchin, D.~Barnes, and C.~Roitzheim.
\newblock N{$_\infty$}-operads and associahedra.
\newblock {\em Pacific J. Math.}, 315(2):285--304, 2021.

\bibitem{HillSlicePrimer}
M.~A. Hill.
\newblock The equivariant slice filtration: a primer.
\newblock {\em Homology Homotopy Appl.}, 14(2):143--166, 2012.

\bibitem{HHR}
M.~A. Hill, M.~J. Hopkins, and D.~C. Ravenel.
\newblock On the nonexistence of elements of {K}ervaire invariant one.
\newblock {\em Ann. of Math. (2)}, 184(1):1--262, 2016.

\bibitem{HillYarnall}
M.~A. Hill and C.~Yarnall.
\newblock A new formulation of the equivariant slice filtration with
  applications to {$C_p$}-slices.
\newblock {\em Proc. Amer. Math. Soc.}, 146(8):3605--3614, 2018.

\bibitem{Rubin21}
J.~Rubin.
\newblock Detecting {S}teiner and linear isometries operads.
\newblock {\em Glasg. Math. J.}, 63(2):307--342, 2021.

\bibitem{WilsonSlice}
D.~Wilson.
\newblock On categories of slices.
\newblock arxiv.org: 1711.03472, 2017.

\end{thebibliography}

\end{document}
